% ============================================================================
%  Chapter 2. Background and Related Work
%
%  \input file. Requires booktabs, amsmath, natbib in the parent preamble.
%  Bibliography: accord/thesis/refs.bib
%
%  Positioning chapter. Contains no results. Where a gap identified here is
%  weakened by this thesis's own findings, that is stated in the same section
%  rather than deferred to Chapter 9.
% ============================================================================

\chapter{Background and Related Work}
\label{ch:background}

\section{How This Chapter Is Organised}
\label{sec:bg-overview}

This chapter has two jobs. The first is conventional: to situate cross-note
medication contradiction detection in the literature it draws on, which spans
five loosely connected clusters, and to identify the gaps the project was
designed to fill. The second is less conventional and is the reason the chapter
appears in this form. Three bodies of work that the project needed, and did not
consult until after its results were in, are developed here at length: the class
imbalance literature, the reporting standards for clinical prediction models,
and the behaviour of agreement and correlation statistics at low prevalence.
Their absence from the original design is not incidental to the failures
documented in Chapters~\ref{ch:weak} through~\ref{ch:eval}. It is close to
sufficient to explain them.

Sections~\ref{sec:bg-nli} through~\ref{sec:bg-safety} cover the five clusters.
Section~\ref{sec:bg-gaps} states the six gaps the project set out to address,
and marks two of them as weakened by this thesis's own findings.
Sections~\ref{sec:bg-imbalance} through~\ref{sec:bg-metrics} cover the three
required additions. Section~\ref{sec:bg-position} states the position that
survives all of it.

\section{Cluster 1: Clinical Natural Language Inference}
\label{sec:bg-nli}

\subsection{The paradigm}

Natural language inference is the task of classifying the relation between a
premise and a hypothesis as entailment, contradiction, or neutral. Large
general-domain corpora \citep{bowman2015snli,williams2018mnli} established the
task as a standard probe of sentence-level semantic competence, and MedNLI
\citep{romanov2018mednli} transferred it to the clinical domain by having
clinicians write hypotheses against premises drawn from MIMIC-III patient
records. MedNLI is the closest well-established benchmark to the problem in this
thesis and Chapter~\ref{ch:eval} uses it as a component sanity check.
Subsequent clinical inference resources have moved toward multi-evidence and
numerical reasoning over structured clinical documents \citep{jullien2023nli4ct},
which is closer to the practical setting but still operates within a document
rather than across a pair of documents authored by different roles.

\subsection{The assumption these resources share}

Every one of these datasets presents the model with an ordered pair in which the
premise is stipulated to be true and the hypothesis is the object of judgement.
The relation is defined semantically and is indifferent to who wrote either
side. This is exactly right for the task as posed, and it is the assumption that
the construct developed in Chapter~\ref{ch:data} as \emph{clinical note
priority} was meant to relax: in a real medication disagreement the two sides
are not interchangeable, because one of them is an order that will be acted on
and the other is a record of what happened.

\subsection{What the literature already establishes about how such models
succeed}

A substantial body of work shows that natural language inference models
frequently reach high benchmark accuracy by exploiting properties of the
hypothesis alone \citep{poliak2018hypothesis,gururangan2018artifacts} or
lexical-overlap heuristics that fail under controlled perturbation
\citep{mccoy2019hans}. The general phenomenon has been characterised as shortcut
learning \citep{geirhos2020shortcut}, and its clinical analogue is documented in
imaging, where models key on site-identifying artifacts and collapse across
institutions \citep{zech2018generalization}. Chapter~\ref{ch:arch} adds a
specific instance: an auxiliary head whose target is derivable from the model's
own input features, which reads as a $0.96$ $F_1$ in domain and collapses to
$0.065$ and $0.060$ on two transfer cohorts.

\paragraph{What remains unresolved.} None of this literature addresses the
asymmetric, cross-document, rare-positive setting. The artifacts work concerns
within-pair leakage in balanced benchmarks; the shortcut work concerns
distribution shift. Neither supplies a diagnostic for a pipeline whose
supervision marginal is inverted relative to the truth, which is the failure
mode documented in Chapter~\ref{ch:weak}.

\section{Cluster 2: Contradiction and Inconsistency Detection}
\label{sec:bg-contradiction}

\subsection{The paradigm}

Contradiction detection as a task distinct from general inference was posed by
\citet{demarneffe2008contradictions}, who supplied a typology (negation,
numeric, antonymy, structural, and others) and the observation that contradiction
is rarer and less well-served by entailment-oriented systems than the symmetric
framing suggests. In the biomedical domain, \citet{alamri2016corpus} constructed
a corpus of potentially contradictory claims from cardiovascular research
abstracts, which is the closest existing resource to a clinical contradiction
benchmark and is nonetheless about published claims rather than about care
documentation.

More recent work has moved to document scale. \citet{hsu2021wikicontradiction}
detect self-contradicting Wikipedia articles by learning contradiction-aware
comparisons over all sentence pairs within an article, and
\citet{li2024contradoc} construct a human-annotated dataset of self-contradictions
in long documents and evaluate large language models on it, finding that the
strongest model is still unreliable on contradictions requiring context and
nuance.

\subsection{Three properties that make the clinical cross-note case different}

First, the two sides are separate documents with different authors and different
institutional standing, so the relation is asymmetric in a way that intra-document
self-contradiction is not. Second, the positive class is rare enough that the
detection problem is inseparable from the retrieval problem: no system can
examine every cross-note medication pair in a cohort, so a candidate generator
sits in front of the model and bounds its recall absolutely
(Chapter~\ref{ch:candidates}). Third, and least anticipated by the literature,
the most common contradiction type in the confirmed positives is \emph{omission}:
seven of the twelve pharmacist-confirmed contradictions in
Chapter~\ref{ch:eval} are a drug present in one note and absent from the other.
Omission is not a relation between two sentences. It is a relation between a
sentence and the absence of one, and no pairwise formulation in this cluster
represents it naturally.

\paragraph{What remains unresolved.} There is no cross-note clinical
contradiction benchmark, no shared task, and no published detection system for
the medication case. That absence is what motivated this project, and
Section~\ref{sec:bg-gaps} states it as G1 together with the sense in which this
thesis weakens it.

\section{Cluster 3: Clinical and Biomedical Language Model Pretraining}
\label{sec:bg-pretraining}

\subsection{The paradigm}

Transformer encoders pretrained on general text \citep{devlin2019bert} lose
substantial accuracy on biomedical and clinical text, and the field's response
has been domain-specific pretraining. The two strategies are continued
pretraining from a general checkpoint and pretraining from scratch on in-domain
text. BioBERT \citep{lee2020biobert} and SciBERT \citep{beltagy2019scibert}
represent the first, on PubMed and on scientific text respectively;
BioMedBERT \citep{gu2021biomedbert} represents the second and reports that
from-scratch domain pretraining, with an in-domain vocabulary, outperforms
continued pretraining across a biomedical benchmark suite. For clinical rather
than biomedical text, Bio\_ClinicalBERT \citep{alsentzer2019clinicalbert}
continues pretraining on MIMIC-III notes, and a parallel line
\citep{huang2019clinicalbert} targets clinical notes for downstream prediction.
Scale has since been pushed considerably further on institutional clinical
corpora \citep{yang2022gatortron}.

This thesis uses BioMedBERT as its backbone, adapted with LoRA
\citep{hu2022lora}. That choice is defensible on published benchmark
performance and is not among the findings under dispute.

\subsection{Two assumptions in this cluster that this thesis strains}

\paragraph{That the head of a note is representative of it.} These models are
almost universally applied with head truncation to a fixed token budget. For
MIMIC-III discharge summaries the head is the administrative header block, and
de-identification surrogates \citep{neamatullah2008deid} are densest exactly
there. Chapter~\ref{ch:data} records this and Chapter~\ref{ch:arch} identifies
it as one of two visible mechanisms behind the collapse of self-supervised
pretraining to the constant-prediction solution. The assumption is not stated
anywhere in the pretraining literature because it is not usually load-bearing;
here it is.

\paragraph{That parameter-efficient adaptation is approximately loss-free.}
LoRA is widely described as recovering full fine-tuning performance at a small
fraction of the trainable parameters \citep{hu2022lora}. Chapter~\ref{ch:eval}
reports that on MedNLI, at $r = 16$ and a shorter schedule, the gap is roughly
eleven accuracy points, $83.90\%$ against $72.93\%$. The two configurations
differ in both rank and schedule and so do not isolate either, which is recorded
as a [TODO] there. The relevant point for this chapter is that the assumption
was inherited rather than checked.

\paragraph{What remains unresolved.} No clinical pretraining objective in this
cluster targets \emph{cross-document} consistency. Chapter~\ref{ch:arch}
describes an attempt at one, and its collapse.

\section{Cluster 4: Weak Supervision and Data Programming}
\label{sec:bg-weak}

\subsection{The paradigm}

Data programming \citep{ratner2016dataprogramming} and its system realisation
Snorkel \citep{ratner2017snorkel} replace hand annotation with multiple noisy
labelling functions and a generative model that estimates their accuracies and
correlations in order to produce probabilistic training labels. The approach has
a substantial clinical track record: ontology-driven weak supervision matches
hand-labelled training data for clinical entity classification
\citep{fries2021trove}, and cross-modal data programming enables rapid model
development in medical imaging by writing rules over an auxiliary text modality
\citep{dunnmon2020crossmodal}. The paradigm and its variants are surveyed in
\citet{zhang2022wssurvey}.

Since 2023 the dominant cheap labeller in practice has been an instruction-tuned
language model, on the strength of results reporting that such models match or
exceed crowd annotators on several classification tasks
\citep{gilardi2023chatgpt}. \citet{pangakis2023validation} respond that any
automated annotation procedure must be validated against human labels on the
same items before its output is used.

\subsection{The assumption this cluster rests on, and where it breaks}

The theoretical guarantees in this cluster require \emph{multiple} sources whose
errors are conditionally independent given the true label. Under that condition
disagreement is informative, accuracies are identifiable, and aggregation
recovers a better label than any single source. The condition is doing real
work, and it is the condition Chapter~\ref{ch:weak} violates in the most direct
way available: there is one source, and its error is not noise but a systematic
prior shift, a positive rate of $84.9\%$ against a measured base rate of
$4.0\%$. Aggregation over one source is the identity function, and no
noise-robust loss removes a shifted marginal.

A second, quieter assumption concerns confidence. Language models emit
confidence values readily and those values are not calibrated probabilities
without measurement \citep{guo2017calibration}, are strongly sensitive to prompt
construction and class priors \citep{zhao2021calibrate}, and can be degenerate.
Chapter~\ref{ch:weak} reports the degenerate case: the confidence field takes
exactly two values across all $1{,}679$ training rows and all $299$ audit rows,
and is a deterministic function of the label.

A third concerns the ground truth against which everything is measured. Label
errors are pervasive even in curated benchmark test sets and materially
destabilise reported rankings \citep{northcutt2021labelerrors}; the situation
here is worse, because the gold labels were collected against evidence that
Chapter~\ref{ch:candidates} shows was frequently the wrong evidence.

\paragraph{What remains unresolved.} The literature contains no protocol for
applying language model weak supervision to a rare clinical class. The
constructive contribution of Chapter~\ref{ch:weak} is a candidate protocol,
which is deliberately unoriginal in its parts, and whose omission is what cost
this project its supervised half.

\section{Cluster 5: Medication Safety and Adverse Drug Event Detection}
\label{sec:bg-safety}

\subsection{The clinical problem}

Adverse drug events are a first-order inpatient safety problem
\citep{bates1995ade}, and the transition from inpatient to outpatient care
concentrates the risk. Unintended medication discrepancies are present at a
large share of admissions \citep{cornish2005discrepancies} and discharges
\citep{wong2008reconciliation}, with omission of a regularly used medication the
single most common category. Errors traceable to discontinuity of care affect a
substantial fraction of discharged patients \citep{moore2003discontinuity}, and
adverse events in the weeks after discharge are common, frequently drug-related,
and frequently preventable \citep{forster2003adverse}. The information transfer
intended to prevent them is unreliable in both content and timing
\citep{kripalani2007deficits}, and discharge summary timeliness itself is
associated with patient outcomes \citep{li2013timeliness}.

Chapter~\ref{ch:eval} reports that omission accounts for seven of the twelve
pharmacist-confirmed contradictions, which is consistent with this literature and
is the reason Chapter~\ref{ch:eval} regards the architecture as structurally
mismatched to its own target population.

\subsection{The NLP problem}

Medication information extraction from clinical text is a mature task with
established shared tasks and strong published systems
\citep{uzuner2010medication,henry2020n2c2}, mature negation and context handling
\citep{chapman2001negex}, and mature open tooling
\citep{neumann2019scispacy,eyre2021medspacy}. Chapter~\ref{ch:data} records that
this project used none of it, for an environment reason rather than an
evidential one, and that the resulting dictionary extractor was never validated
against a human-annotated mention set. That is the most consequential missing
measurement in the data layer, because it means no downstream error rate can be
decomposed into extraction error and model error.

\subsection{The deployment problem}

The clinical decision support literature is unambiguous about what happens to a
low-precision alert. Override rates for drug safety alerts in computerised order
entry run between $49\%$ and $96\%$ \citep{vandersijs2006overriding}, and
alerting burden and work complexity degrade response
\citep{ancker2017alertfatigue}. Chapter~\ref{ch:eval}'s ethics statement draws
the direct consequence: a detector at $91.7\%$ recall and $4.7\%$ precision
produces roughly twenty false alarms per true contradiction and should not be
deployed in any form.

\paragraph{What remains unresolved.} This literature establishes the problem's
importance and the cost of a bad solution. It does not supply an automated
cross-note detection method, which is the gap the project targeted.

\section{The Six Gaps, and Which of Them This Thesis Weakens}
\label{sec:bg-gaps}

The project was designed against six gaps. They are restated here as they were
posed, with an explicit note where the thesis's own findings weaken the framing.

\begin{description}
\item[G1. No cross-note clinical contradiction system.] Contradiction detection
      exists within documents and between research claims; nothing addresses
      disagreement between two care documents of different types about the same
      medication in the same admission.

      \textit{Weakened.} The gap is real as a statement about the literature.
      But this thesis's own results show that the framing of the gap as a
      \emph{modelling} gap was wrong. Chapter~\ref{ch:candidates} establishes
      that the binding constraint is the candidate generator, whose pool sits at
      $4.0\%$ precision, and Chapter~\ref{ch:eval} establishes that no
      downstream model repairs that. The unfilled gap is a retrieval and
      evidence-attribution gap, not an inference gap, and stating it as the
      latter is what led to a project that built an inference model.

\item[G2. Note-type authority is not represented in clinical inference.] All
      clinical inference resources treat the premise-hypothesis relation
      symmetrically with respect to authorship. Institutional note priority is a
      real feature of clinical documentation and is not modelled.

      \textit{Weakened.} Chapter~\ref{ch:arch} falsifies the claim that this
      project's mechanism learned note priority, four independent ways, and shows
      that constraining the model to symmetry \emph{improves} both reported
      metrics. That is not evidence that G2 is a false gap; it is evidence that
      this thesis did not test it, because the parameter was initialised at the
      values later reported as discovered and sits on a branch a subsequent
      normalisation largely erases. G2 remains open, and
      Chapter~\ref{ch:discussion} specifies what a correct test would require.

\item[G3. No guideline-grounded severity ranking.] Confirmed contradictions
      differ enormously in danger, and no system ranks them using evidence from
      published practice guidelines.

      \textit{Untested rather than addressed.} Chapter~\ref{ch:arch} establishes
      that the component intended to fill this gap has no guideline corpus and
      receives no gradient. G3 is exactly as open as it was.

\item[G4. Temporal context is unencoded.] Cross-note reasoning should be
      sensitive to the order of and interval between the two documents.

      \textit{Unaddressable in this resource.} \texttt{CHARTTIME} is null for
      $100\%$ of MIMIC-III discharge summaries (Chapter~\ref{ch:data},
      Section~\ref{sec:charttime}). The gap stands; MIMIC-III cannot be used to
      close it at sub-day resolution.

\item[G5. No weak supervision protocol for contradiction labelling.] Data
      programming has clinical precedent for entity and image tasks but not for
      rare pairwise judgements.

      \textit{Confirmed, and given a candidate answer.} Chapter~\ref{ch:weak}
      documents why the naive application fails and states a protocol of six
      pre-training checks. The protocol is not novel in its parts. Its value is
      that it is stated as a stop condition rather than as advice.

\item[G6. Self-supervised clinical coherence pretraining has not been
      attempted.] No clinical pretraining objective targets whether two
      documents from the same episode are mutually consistent.

      \textit{Attempted and collapsed.} Chapter~\ref{ch:arch} reports
      convergence to the constant-prediction solution at $\ln 2$ with
      development accuracy exactly $0.5000$, identically under a $37.9\times$
      change in trainable capacity, and withdraws the explanation originally
      offered for it as untested.
\end{description}

\section{Required Addition 1: Class Imbalance}
\label{sec:bg-imbalance}

\subsection{Why this section exists}

The project treated its $4\%$ positive rate as a difficulty to be managed by the
model rather than as a property that determines what any measurement of the
model can mean. The class imbalance literature is large, mature, and was not
consulted. This section summarises what it offers, and then states the more
important point: none of it addresses the problem this thesis actually ran into,
which is an evaluation problem rather than a training problem.

\subsection{Four families of response}

\paragraph{Resampling.} Undersample the majority, oversample the minority, or
synthesise minority examples by interpolation \citep{chawla2002smote}. The
systematic study of \citet{japkowicz2002imbalance} shows that the difficulty of
an imbalanced problem depends jointly on imbalance ratio, concept complexity,
and training set size rather than on imbalance alone, which already predicts
that a $4\%$ rate over $1{,}679$ examples of a semantically subtle relation is
an unfavourable corner of that space. The broad survey is
\citet{he2009imbalanced}.

\paragraph{Cost-sensitive learning.} Rather than change the data, change the
decision rule. \citet{elkan2001cost} gives the foundational result: for a
two-class problem the optimal decision threshold is determined by the
misclassification cost ratio, and rebalancing the training distribution is
equivalent to a particular reweighting rather than a distinct method. This is
the correct frame for a clinical triage system, where the cost ratio between a
missed dangerous contradiction and a false alarm is very far from one, and it is
also the frame this project never made explicit.

\paragraph{Loss reshaping.} Down-weight easy examples \citep{lin2017focal},
weight classes by effective sample count \citep{cui2019classbalanced}, enforce
label-distribution-aware margins \citep{cao2019ldam}, or adjust logits by label
frequency either post hoc or in the loss \citep{menon2021logitadjust}. These
methods are principled and effective when the model has learned a conditional
worth reshaping.

\paragraph{Positive-unlabelled learning.} When negatives are not reliably
observed, treat the problem as learning from positive and unlabelled data
\citep{bekker2020pu}. This is arguably the closest formal match to the setting
in Chapter~\ref{ch:candidates}, where a rule fires and the non-firing remainder
is unlabelled rather than known-negative, and it is not attempted anywhere in
this thesis.

\subsection{Why none of this addresses the evaluation problem}

Two results in this literature bear directly on the thesis and should be read
together.

First, \citet{vandengoorbergh2022harm} show, in a clinical prediction setting,
that imbalance corrections including undersampling, oversampling, and SMOTE
produce severe miscalibration, systematically overestimating the probability of
minority-class membership, without improving discrimination as measured by the
area under the ROC curve. Whatever gain in sensitivity and specificity these
methods deliver is obtainable by shifting the decision threshold instead. The
implication is that imbalance corrections move the operating point and do not
create ranking information.

Second, that is precisely the quantity this thesis has none of.
Chapter~\ref{ch:eval} reports AUROC $0.5145$ with a Mann-Whitney $p$ of $0.867$,
and $0.5078$ with $p = 0.953$ on the strata free of the exemplar defect. A
threshold shift, a resampling scheme, a focal loss, or a logit adjustment all
operate on a score ordering. When the ordering is uninformative, every one of
them is a relabelling. This is why Chapter~\ref{ch:weak} can state that prior
correction would not have sufficed \emph{here} while declining to claim that
prior correction is generally ineffective.

The deeper point, which Chapter~\ref{ch:metrics} develops in full, is that the
imbalance literature is overwhelmingly about \emph{fitting} under skew and only
incidentally about \emph{selecting} under skew. The question of which system a
summary statistic prefers when the positive class is $4.74\%$ of an evaluation
set, and what exchange rate between missed positives and avoided false alarms
that preference encodes, is not answered by any method in this section. The
tools that do answer it are threshold and decision analysis
\citep{vickers2006dca} and an explicit statement of the cost ratio
\citep{elkan2001cost}, and neither is standard practice in the NLP literature
this project was written for.

\section{Required Addition 2: Reporting Standards for Clinical Prediction Models}
\label{sec:bg-tripod}

A clinical contradiction detector is a clinical prediction model, and clinical
prediction models have a reporting standard. TRIPOD+AI
\citep{collins2024tripodai} updates the earlier TRIPOD guidance to cover models
developed with regression or machine learning, and applies to development,
validation, and combined designs. It is adopted here as the reporting frame for
the thesis, and the adoption is not decorative: several of the failures
documented in Chapters~\ref{ch:candidates} through~\ref{ch:eval} correspond to
checklist items that were not satisfied.

Four correspondences are worth naming explicitly.

\paragraph{Participants and data sources.} TRIPOD+AI requires the eligibility
criteria, the data source, and the flow of participants through the study to be
reported. Chapter~\ref{ch:data} supplies this as a cohort funnel, including the
six note-role assignment attempts and the retreat from three roles to two, and
including a filter that was computed and printed but never applied to the set
actually joined against.

\paragraph{Outcome definition and blinding.} The outcome here is a pharmacist's
binary contradiction judgement, and the checklist requires that how it was
determined, and by whom, be stated. Chapter~\ref{ch:candidates} states it: one
annotator, two passes, a two-week washout, no established inter-rater
reliability, and a previously reported $\kappa = 1.000$ retracted as circular
because its two ``instruments'' stood in an identity relation.

\paragraph{Sample size.} TRIPOD+AI requires an explicit justification of sample
size rather than a convenience sample reported without comment.
Chapter~\ref{ch:discussion} states the gap between the twelve positives
available and the prospectively stated requirement, and reports every interval
at the width that twelve positives permit.

\paragraph{Model specification and performance measures.} The checklist requires
that the model actually fitted be described, and that the performance measures
be prespecified. Chapter~\ref{ch:arch} is, in effect, a report of what happens
when the first requirement is unmet: three components of the specification are
not present in the code that produced the results, one of them replaced by a
vector of zeros and one by twenty randomly initialised vectors that receive no
gradient. Chapter~\ref{ch:metrics} is a report of what happens when the second
is unmet, which is that the metric ends up chosen after the results are known.

\paragraph{Positioning.} Nothing in this thesis argues that TRIPOD+AI compliance
would have produced a working system. It argues something narrower and, for a
reader deciding what to do next, more useful: the checklist items above are
mechanical, they are checkable before a results table exists, and satisfying
four of them would have surfaced four of this project's failures at the point
where they were still cheap to fix.

\section{Required Addition 3: Evaluation Metrics Under Low Prevalence}
\label{sec:bg-metrics}

\subsection{The measures in question}

Three statistics recur in this thesis and each behaves in a way that is
counterintuitive when the positive class is rare.

\paragraph{Matthews correlation.} Introduced for the assessment of secondary
structure prediction \citep{matthews1975mcc}, MCC is the Pearson correlation
between the binary prediction and the binary label, computed from all four cells
of the confusion matrix. \citet{chicco2020mcc} argue that it should be preferred
to accuracy and $F_1$ under imbalance, on the grounds that it is high only when
all four cells are favourable, and \citet{powers2011evaluation} situates it in a
larger family alongside informedness and markedness. Chapter~\ref{ch:eval}
reports MCC for every compared system for exactly this reason, and
Chapter~\ref{ch:metrics} shows what it selects for at $P = 12$, $Q = 241$.

\paragraph{Cohen's kappa and its prevalence-adjusted variant.} Kappa
\citep{cohen1960kappa} corrects observed agreement for agreement expected by
chance under the observed marginals, which makes it sensitive to those
marginals in ways that are widely regarded as paradoxical:
\citet{feinstein1990paradoxes} show that a high observed agreement can be driven
to a low kappa by marginal imbalance, and \citet{cicchetti1990resolving} analyse
the resolution. \citet{byrt1993pabak} respond with PABAK, the
prevalence-adjusted bias-adjusted kappa, which is a simple affine function of
observed agreement, $\mathrm{PABAK} = 2p_o - 1$, and therefore removes the
marginal sensitivity by removing the marginal information.
\citet{delgado2019kappa} argue on the strength of these behaviours that kappa
should be avoided as a classification performance measure in favour of MCC.

\paragraph{Balanced accuracy.} The unweighted mean of sensitivity and
specificity \citep{brodersen2010balanced}, which is exactly the quantity a
prevalence-blind summary would want and, as Chapter~\ref{ch:metrics} shows in
closed form, exactly the quantity that encodes an implicit exchange rate of
$Q/P$ true negatives per true positive.

\subsection{Why average precision is reported alongside AUROC}

Under heavy negative skew the ROC curve is dominated by the majority class and
can look acceptable while the precision-recall behaviour is at the no-skill line
\citep{saito2015prcurve}. Chapter~\ref{ch:eval} therefore reports both, and the
two agree: AUROC $0.5145$ and average precision $0.0579 \pm 0.0034$
against a base rate of $0.040$.

\subsection{The gap this literature leaves}

Each of these papers establishes that a metric misbehaves under imbalance, and
recommends a different metric. What none of them does is state, for a given
evaluation set, the \emph{exchange rate} that a chosen metric imposes between
false negatives and false positives, in a form a reviewer can check in one line.
Chapter~\ref{ch:metrics} supplies that for this evaluation set, and makes the
recommendation prescriptive rather than descriptive: it does not claim the field
selects models this way, because that evidence does not exist and was correctly
demanded when the claim was previously made. It claims that if a paper selects
by these statistics at this prevalence, the exchange rate it is accepting should
be stated in the paper.

\section{Position}
\label{sec:bg-position}

The five clusters supply, respectively: a task formulation that is symmetric
where the clinical relation is not; a contradiction literature that is
intra-document or inter-claim rather than inter-document and does not represent
omission; a pretraining literature whose truncation and adaptation assumptions
are load-bearing here and were inherited unchecked; a weak supervision
literature whose guarantees require multiple conditionally independent sources
and were applied to one; and a clinical safety literature that establishes the
problem's importance and the cost of a low-precision solution without supplying a
method.

The three required additions supply what the project needed and did not have: a
frame in which a $4\%$ positive rate is a property of the measurement problem
and not only of the fitting problem; a checklist whose mechanical items would
have surfaced four of this thesis's failures before they reached a results
table; and a metrics literature that establishes these statistics misbehave under
imbalance without stating what they trade.

The position this thesis takes is therefore not that it fills G1 through G6.
Two of the six are weakened as framings by its own findings, one was never
tested, one is unaddressable in this resource, and two are answered negatively.
The position is that the composed pipeline these gaps motivate fails at every
stage, that each failure has an identifiable mechanism and a cheap detector, and
that the composition, not any single component, is what the field should
re-examine.
